
\begin{table}[!htbp]
\centering
\caption{Compared methods and runtime deployment variants.}
\label{tab_compared_methods}
\scriptsize
\setlength{\tabcolsep}{3.4pt}
\renewcommand{\arraystretch}{1.08}
\begin{tabular}{@{}>{\raggedright\arraybackslash}p{0.16\textwidth}>{\raggedright\arraybackslash}p{0.20\textwidth}>{\raggedright\arraybackslash}p{0.25\textwidth}>{\raggedright\arraybackslash}p{0.30\textwidth}@{}}
\toprule
Method & Training or reward & Runtime operator & Role in protocol \\
\midrule
PPO backbone & No risk reward & None & Fixed policy backbone and lower-bound comparator \\
Risk-only & Risk reward & None & Tests whether low cost is stationary avoidance \\
No-action guard & Risk reward & Disabled & Clean reward-only versus runtime-intervention control \\
RSS/TTC filter & No risk reward & External TTC/RSS rule & Classical runtime-filter comparator \\
PPO-Lagrangian & Lagrangian penalty & None & Training-time constrained baseline under this protocol \\
Guard & Curriculum PPO, no risk reward & TTC hard/soft cap + speed guard & Runtime deployment variant with action intervention \\
Shield & No-curriculum PPO, no risk reward & Same TTC hard/soft cap + speed guard & Runtime deployment variant isolating the execution-side operator \\
Gated-risk & Curriculum PPO with gated risk & Same TTC hard/soft cap + speed guard & Adds motion/risk-gated reward to reduce stationary exploitation \\
\bottomrule
\end{tabular}
\begin{tablenotes}
\item Guard and Shield share the same execution-side operator but differ in training context; they are deployment variants rather than a clean single-factor ablation. The clean intervention contrast is Risk-only/No-action guard versus rows with runtime action intervention.
\end{tablenotes}
\end{table}
